\section{Conventions and probability notation}\label{sec:preliminaries}

After stating \Cref{thm:main}, we use unit-period coordinates
$y=x/(2\pi)$ and normalized length $dy$ on $\mathbb R/\mathbb Z$.
Characters are
\[
 \chi_u(y)=e^{2\pi iuy},\qquad u\in\mathbb Z.
\]
For real $u$ the same exponential notation is used on real intervals;
only integer frequencies are regarded as characters of the circle.
Periodic functions are extended to the real line when restricted to
spatial intervals. All conclusions are unchanged by modifications on
null sets.

All logarithms are natural. Write
$\log^-t=\max\{0,-\log t\}$ for $t>0$. The parameter $d$ is sufficiently
large; its lower threshold may depend on fixed auxiliary constants
chosen earlier. Numerical constants denoted by $c,C$ can change from
line to line. Any other dependence is displayed or explicitly stated.
No constant depends on the spatial depth or on a frequency cutoff.
The notation $\operatorname{poly}(x_1,\ldots,x_k)$ denotes a polynomial
upper bound with fixed nonnegative coefficients and fixed finite
exponents. Dyadic levels of a positive parameter are integer powers
of two. Integer roundings that matter are stated explicitly.

For a probability law $P$ on a finite set, its Shannon entropy is
\[
 \mathrm H(P)=-\sum_x P(x)\log P(x),
\]
with $0\log0=0$. Conditional entropy is averaged over the conditioning
law. For general probability measures, relative entropy is
\[
 D(P\|Q)=\int\log\!\left(\frac{dP}{dQ}\right)dP
\]
when $P\ll Q$, with value $+\infty$ otherwise. The same convention applies
when the integral diverges. Mutual information is
$\mathrm I(X;Y)=D(P_{X,Y}\|P_X\otimes P_Y)$, and conditional mutual
information is the corresponding average over the conditioning
variable. A random variable may stand for its law inside $D$ or a
distance between measures.

We use nonnegativity, the chain rules, and contraction of relative
entropy under measurable maps. For background on these quantities and
their chain rules, see \cite[Chapter~2]{CoverThomas2006}. With
\[
 \|P-Q\|_{\mathrm{TV}}=\sup_A|P(A)-Q(A)|,
\]
Pinsker's inequality takes the form
$\|P-Q\|_{\mathrm{TV}}\le\sqrt{D(P\|Q)/2}$. For probability densities
this is one half of their $L^1$ distance. Numerical constants absorb
the harmless choice between total variation and $L^1$ conventions when
an estimate is stated with an unspecified constant.

In the entropy compression argument, all Shannon entropies concern
finite-valued variables. For a finite set $S\subset\mathbb Z$, $nS$
is the $n$-fold sumset with repetitions. Later, $[D]$ denotes the signed
span of a finite list $D$:
\[
 [D]=\left\{\sum_{v\in D}\varepsilon_v v:
                    \varepsilon_v\in\{-1,0,1\}\right\}.
\]
Lists retain their order of insertion; repeated coordinates, if
present, are treated as separately indexed entries.

The dyadic construction uses two probability laws. The law
$\mathbf U$ follows a uniformly chosen spatial output point and
therefore has fair binary bits. The true input law $\mathbf P$ is
weighted by the current nonnegative input. Additional noisy labels are
observations conditional on the input point. The deterministic
processing choices depend on the spatial and membership state, not on
the realized label noise. These distinctions will be used explicitly
in \Cref{sec:procedure}; they are essential when converting information
under the input law into estimates on output paths.
